\documentclass[10pt,openany]{book}
\usepackage[T2A]{fontenc}
\usepackage[utf8]{inputenc}
\usepackage[russian]{babel}
\usepackage{amsmath,amssymb,amsthm,mathtools}
\usepackage{geometry}
% формат книги: 60x90/16 = 145x215 мм, полоса набора 6 1/4 кв. (113 мм)
\geometry{paperwidth=145mm,paperheight=215mm,textwidth=113mm,textheight=178mm,
  inner=14mm,top=20mm,headheight=12pt,headsep=4mm,footskip=8mm,marginparwidth=11.5mm,marginparsep=2.5mm}
\usepackage{enumitem}
\usepackage{tikz}
\usetikzlibrary{patterns,calc,arrows.meta,positioning}
\usepackage[most]{tcolorbox}
\usepackage{xcolor}
\usepackage{marginnote}
\usepackage{booktabs}
\usepackage[hidelinks]{hyperref}

% --- обозначения ---
\newcommand{\ov}[1]{\overline{#1}}
\renewcommand{\emptyset}{\varnothing}
\newcommand{\R}{\mathbb R}
\newcommand{\N}{\mathbb N}
\newcommand{\Z}{\mathbb Z}
\newcommand{\Q}{\mathbb Q}
\newcommand{\M}{\mathfrak M}
\newcommand{\maj}{\operatorname{maj}}
\newcommand{\sd}{\mathbin{\triangle}}
\newcommand{\cupl}{\mathbin{\cup}}
% номер страницы оригинала на полях
\newcommand{\bookpage}[1]{\leavevmode\marginnote{\footnotesize\textcolor{gray}{[с.~#1]}}\ignorespaces}
% подпись к рисункам в стиле книги
\newcommand{\figcap}[1]{\par\smallskip\centering\small #1\par}
% --- добавленный материал ---
\definecolor{solcol}{RGB}{20,80,150}
\newtcolorbox{solution}[1][]{breakable,enhanced,colback=solcol!4,colframe=solcol!60,
  boxrule=0.5pt,arc=1pt,left=6pt,right=6pt,top=4pt,bottom=4pt,
  title={Решение упражнения #1},fonttitle=\bfseries\small,coltitle=white,colbacktitle=solcol!75}
\newtcolorbox{remarkbox}[1][]{breakable,enhanced,colback=orange!5,colframe=orange!60!black,
  boxrule=0.4pt,arc=1pt,title={#1},fonttitle=\bfseries\small}
% исправление опечатки: отметка в тексте
\newcommand{\fixed}[1]{\footnote{\textit{Исправлено:} #1}}
\theoremstyle{plain}
\newtheorem*{lemma*}{Лемма}
\newtheorem*{theorem*}{Теорема}
\newtheorem*{prop*}{Утверждение}
\newcommand{\exitem}[1]{\par\smallskip\noindent\textbf{#1}~}

\emergencystretch=3em
\tolerance=1500
\usepackage{fancyhdr}
\pagestyle{fancy}\fancyhf{}
\fancyhead[LE,RO]{\thepage}\fancyhead[RE]{\small\itshape Рвачев. Решения упражнений}\fancyhead[LO]{\small\itshape\nouppercase{\rightmark}}
\fancypagestyle{plain}{\fancyhf{}\fancyfoot[C]{\thepage}\renewcommand{\headrulewidth}{0pt}}

\title{В.\,Л.~Рвачев\\[4pt]
\Large Теория $R$-функций и некоторые ее приложения\\[8pt]
\large Решения упражнений}
\author{}
\date{}

\begin{document}
\maketitle
\tableofcontents

\chapter*{Глава 1. Элементы дискретной математики}
\addcontentsline{toc}{chapter}{Глава 1. Элементы дискретной математики}
\markboth{Глава 1}{Глава 1}
\section*{\S\,1. Множества и логические операции над ними}
\addcontentsline{toc}{section}{\S\,1. Множества и логические операции над ними}
\markright{\S\,1. Множества и логические операции над ними}
\subsection*{Условия}
1. Доказать тождества
\begin{equation}\tag{1.4}
\begin{gathered}
A\sd B\equiv(A\cup B)\setminus(A\cap B);\\
A\cup B\equiv(A\sd B)\sd(A\cap B);\qquad A\setminus B\equiv A\sd(A\cap B);\\
(A\cap B)\cup B\equiv B;\qquad (A\cup B)\cap B\equiv B;\\
(A\cap B)\cup(A\cap\ov B)\equiv A;\qquad (A\cup B)\cap(A\cup\ov B)\equiv A;\\
(\ov A\cap B)\cup(A\cap\ov B)\equiv(A\cup B)\cap(\ov A\cup\ov B).
\end{gathered}
\end{equation}
(Для доказательства использовать формулы (1.1), (1.2) и свойства
$1^\circ$--$19^\circ$.)

2. Показать, что симметрическая разность "--- это коммутативная
и ассоциативная операции.

\textsc{Решение.} В соответствии с формулой (1.2) и свойствами
$1^\circ$--$19^\circ$ получаем
\begin{align*}
A\sd(B\sd C)&\equiv[A\cap(\ov{B\sd C})]\cup[\ov A\cap(B\sd C)]\equiv\\
&\equiv\{A\cap\ov{[(B\cap\ov C)\cup(\ov B\cap C)]}\}\cup\{\ov A\cap[(B\cap\ov C)\cup(\ov B\cap C)]\}\equiv\\
&\equiv\{A\cap[\ov{(B\cap\ov C)}\cap\ov{(\ov B\cap C)}]\}\cup\{\ov A\cap[(B\cap\ov C)\cup(\ov B\cap C)]\}\equiv\\
&\equiv\{A\cap[(\ov B\cup C)\cap(B\cup\ov C)]\}\cup\{(\ov A\cap B\cap\ov C)\cup(\ov A\cap\ov B\cap C)\}\equiv\\
&\equiv\{[A\cap(\ov B\cup C)]\cap[A\cap(B\cup\ov C)]\}\cup(\ov A\cap B\cap\ov C)\cup{}\\
&\qquad\cup(\ov A\cap\ov B\cap C)\equiv\{[(A\cap\ov B)\cup(A\cap C)]\cap[(A\cap B)\cup(A\cap\ov C)]\}\cup{}\\
&\qquad\cup(\ov A\cap B\cap\ov C)\cup(\ov A\cap\ov B\cap C)\equiv\\
&\equiv\{[(A\cap\ov B)\cap(A\cap B)]\cup[(A\cap\ov B)\cap(A\cap\ov C)]\cup[(A\cap C)\cap(A\cap B)]\cup{}\\
&\qquad\cup[(A\cap C)\cap(A\cap\ov C)]\}\cup(\ov A\cap B\cap\ov C)\cup(\ov A\cap\ov B\cap C)\equiv\\
&\equiv\{[(A\cap A)\cap(B\cap\ov B)]\cup[(A\cap A)\cap(\ov B\cap\ov C)]\cup[(A\cap A)\cap(C\cap B)]\cup{}\\
&\qquad\cup[(A\cap A)\cap(C\cap\ov C)]\}\cup(\ov A\cap B\cap\ov C)\cup(\ov A\cap\ov B\cap C)\equiv\\
&\equiv(A\cap\emptyset)\cup(A\cap\ov B\cap\ov C)\cup(A\cap B\cap C)\cup(A\cap\emptyset)\cup{}\\
&\qquad\cup(\ov A\cap B\cap\ov C)\cup(\ov A\cap\ov B\cap C)\equiv\emptyset\cup(A\cap\ov B\cap\ov C)\cup{}\\
&\qquad\cup(A\cap B\cap C)\cup\emptyset\cup(\ov A\cap B\cap\ov C)\cup(\ov A\cap\ov B\cap C)\equiv\\
&\equiv(A\cap\ov B\cap\ov C)\cup(A\cap B\cap C)\cup(\ov A\cap B\cap\ov C)\cup(\ov A\cap\ov B\cap C).
\end{align*}
Поскольку последнее выражение инвариантно относительно циклической
перестановки букв (т.\,е. при замене $A$, $B$, $C$ на $C$, $A$, $B$
соответственно), то оно равно также $C\sd(A\sd B)$. Таким образом,
ассоциативность операции $A\sd B$ доказана. Ее коммутативность
очевидна.\fixed{в оригинале в восьмой строке выкладки стоит скобка
<<$([A\cap A)$>> вместо <<$[(A\cap A)$>>.}

3. Пусть $I=\R$. Показать, что равенство
$(A\cap B)\cup C=C\cup(B\cap C)$ справедливо, если $A=[a,b]$, $B=[c,d]$,
$C=[e,f]$ и $e\le\max(a,c)$; $f\ge\min(b,d)$, но оно не является
тождеством.

4. Показать, что правила де Моргана $10^\circ$ и $11^\circ$ верны для
$n$ ($n\ge2$) компонент, т.\,е.
\begin{equation}\tag{1.5}
\ov{\bigcap_{i=1}^{i=n}A_i}\equiv\bigcup_{i=1}^{i=n}\ov{A_i};\qquad
\ov{\bigcup_{i=1}^{i=n}A_i}\equiv\bigcap_{i=1}^{i=n}\ov{A_i}.
\end{equation}

\subsection*{Решения}
\begin{solution}[1 (\S\,1)]
Будем пользоваться только (1.1), (1.2) и свойствами $1^\circ$--$19^\circ$.
Полезно заранее заметить два следствия.

(а) \emph{Поглощение:} $(A\cap B)\cup B\equiv B$ и $(A\cup B)\cap B\equiv B$.
Действительно,
$(A\cap B)\cup B\overset{14^\circ}{\equiv}(A\cap B)\cup(I\cap B)
\overset{5^\circ}{\equiv}(A\cup I)\cap B\overset{15^\circ}{\equiv}I\cap B\equiv B$;
второе тождество двойственно первому (те же шаги с $6^\circ$, $17^\circ$,
$16^\circ$):
$(A\cup B)\cap B\equiv(A\cup B)\cap(\emptyset\cup B)\equiv(A\cap\emptyset)\cup B\equiv\emptyset\cup B\equiv B$.
Это четвертое и пятое тождества (1.4).

(б) \emph{Склеивание:}
$(A\cap B)\cup(A\cap\ov B)\overset{5^\circ}{\equiv}A\cap(B\cup\ov B)
\overset{13^\circ}{\equiv}A\cap I\equiv A$;
$(A\cup B)\cap(A\cup\ov B)\overset{6^\circ}{\equiv}A\cup(B\cap\ov B)
\overset{12^\circ}{\equiv}A\cup\emptyset\equiv A$. Это шестое и седьмое тождества.

\smallskip
\emph{Последнее тождество.} Раскроем скобки в правой части по $5^\circ$:
\[
(A\cup B)\cap(\ov A\cup\ov B)\equiv(A\cap\ov A)\cup(A\cap\ov B)\cup(B\cap\ov A)\cup(B\cap\ov B)
\equiv(A\cap\ov B)\cup(\ov A\cap B),
\]
так как $A\cap\ov A\equiv B\cap\ov B\equiv\emptyset$. Попутно получена
полезная формула $A\sd B\equiv(A\cup B)\cap(\ov A\cup\ov B)$.

\emph{Первое тождество.} По (1.1) и $10^\circ$:
$(A\cup B)\setminus(A\cap B)\equiv(A\cup B)\cap\ov{A\cap B}\equiv(A\cup B)\cap(\ov A\cup\ov B)$,
а это, как только что показано, $A\sd B$.

\emph{Третье тождество.} По (1.2), $10^\circ$ и $9^\circ$:
\begin{multline*}
A\sd(A\cap B)\equiv\bigl(A\cap\ov{A\cap B}\bigr)\cup\bigl(\ov A\cap A\cap B\bigr)
\equiv\bigl(A\cap(\ov A\cup\ov B)\bigr)\cup\emptyset\equiv{}\\
\equiv(A\cap\ov A)\cup(A\cap\ov B)\equiv A\cap\ov B\equiv A\setminus B.
\end{multline*}

\emph{Второе тождество.} Обозначим $D=A\sd B=(A\cap\ov B)\cup(\ov A\cap B)$,
$P=A\cap B$. Заметим, что $D\cap P\equiv\emptyset$ (каждый член $D$
содержит $\ov A$ или $\ov B$). Тогда $D\cap\ov P\equiv D$ (так как
$D\equiv(D\cap P)\cup(D\cap\ov P)$ по склеиванию) и аналогично
$P\cap\ov D\equiv P$. Поэтому
\[
(A\sd B)\sd(A\cap B)\equiv(D\cap\ov P)\cup(\ov D\cap P)\equiv D\cup P
\equiv(A\cap\ov B)\cup(\ov A\cap B)\cup(A\cap B).
\]
Склеивая $(A\cap\ov B)\cup(A\cap B)\equiv A$ и $(\ov A\cap B)\cup(A\cap B)\equiv B$
(член $A\cap B$ можно повторить по $8^\circ$), получаем $A\cup B$.
\end{solution}

\begin{solution}[2 (\S\,1): дополнение]
Решение в книге приводит $A\sd(B\sd C)$ к виду
$K_{100}\cup K_{111}\cup K_{010}\cup K_{001}$, где
$K_{\alpha\beta\gamma}=A^\alpha\cap B^\beta\cap C^\gamma$ ($X^1=X$,
$X^0=\ov X$). Это ровно те <<кусочки>> диаграммы Эйлера, в которых точка
принадлежит \emph{нечетному} числу множеств. Отсюда сразу видна и
коммутативность, и ассоциативность: для любых $A_1,\dots,A_n$ множество
$A_1\sd\dots\sd A_n$ (при любой расстановке скобок) "--- это множество
точек, принадлежащих нечетному числу из $A_i$. (Индукция: $x\in X\sd Y$
$\iff$ $x$ лежит ровно в одном из $X$, $Y$, т.\,е. число вхождений в $X\sd Y$
имеет ту же четность, что сумма чисел вхождений.) Коммутативность $A\sd B\equiv B\sd A$
следует из (1.2) и $1^\circ$, $2^\circ$.
\end{solution}

\begin{solution}[3 (\S\,1)]
По поглощению $C\cup(B\cap C)\equiv C$, поэтому равенство
$(A\cap B)\cup C=C\cup(B\cap C)$ равносильно $(A\cap B)\cup C=C$, т.\,е.
\[
A\cap B\subset C .
\]
Для отрезков $A\cap B=[\max(a,c),\ \min(b,d)]$ (пусто, если
$\max(a,c)>\min(b,d)$). Если $e\le\max(a,c)$ и $f\ge\min(b,d)$, то
$[\max(a,c),\min(b,d)]\subset[e,f]$, и равенство выполнено (если же
$A\cap B=\emptyset$, оно тем более верно).

Тождеством равенство не является: при $A=B=[0,1]$, $C=[2,3]$ слева
$[0,1]\cup[2,3]$, справа $[2,3]$. Вообще, равенство верно в точности при
$A\cap B\subset C$; условия задачи "--- достаточные (а при
$A\cap B\neq\emptyset$ и необходимые).
\end{solution}

\begin{solution}[4 (\S\,1)]
Индукция по $n$. При $n=2$ это $10^\circ$, $11^\circ$. Пусть (1.5) верно
для $n$. Тогда по ассоциативности, $10^\circ$ и предположению индукции
\[
\ov{\bigcap_{i=1}^{n+1}A_i}\equiv\ov{\Bigl(\bigcap_{i=1}^{n}A_i\Bigr)\cap A_{n+1}}
\equiv\ov{\bigcap_{i=1}^{n}A_i}\cup\ov{A_{n+1}}
\equiv\Bigl(\bigcup_{i=1}^{n}\ov{A_i}\Bigr)\cup\ov{A_{n+1}}\equiv\bigcup_{i=1}^{n+1}\ov{A_i}.
\]
Второе тождество получается так же из $11^\circ$ (или из первого,
примененного к $\ov{A_i}$, с последующим взятием дополнения и законом
$9^\circ$). Можно рассуждать и непосредственно: $x\notin\bigcap A_i$
$\iff$ $x\notin A_i$ хотя бы для одного $i$ $\iff$ $x\in\bigcup\ov{A_i}$.
\end{solution}

\section*{\S\,2. Диаграммы Эйлера}
\addcontentsline{toc}{section}{\S\,2. Диаграммы Эйлера}
\markright{\S\,2. Диаграммы Эйлера}
\subsection*{Условия}
1. Используя диаграммы Эйлера, доказать тождества\fixed{в оригинале последнее тождество напечатано как
$A\cup[(B\setminus\ov A)\sd C]\equiv(A\cap B)\sd C$, которое тождеством
не является (при $A=I$, $B=C=\emptyset$ левая часть равна $I$, правая "---
$\emptyset$). Восстановленный вариант и разбор "--- в решении ниже.}
\[
\begin{gathered}
(A\cup C)\cap B\equiv(A\cap B)\cup(\ov A\cap B\cap C);\\
[A\cap(B\cup C)]\cup(\ov A\cap C)\equiv C\cup(A\cap B);\\
(A\cap C)\cup(B\cap C)\cup(C\setminus A)\equiv C;\qquad A\sd B\equiv(A\cup B)\setminus(A\cap B);\\
[A\cap(B\sd C)]\cup(B\sd C)\equiv B\sd C;\\
A\cap[(B\setminus\ov A)\sd C]\equiv(A\cap B)\sd(A\cap C).
\end{gathered}
\]


2. Выпуклые области $A,B,C\subset\R^2$ разбивают плоскость на восемь частей,
различающихся состояниями принадлежности их точек множествам $A$, $B$ и $C$.
Может ли добавление выпуклой области $D$ привести к разбиению $\R^2$ на 16
подобластей указанного типа? Какое максимальное число таких подобластей
могут образовать $n$ выпуклых областей в $\R^2$?

3. Существуют ли (для всякого $n$) конгруэнтные (т.\,е. получающиеся одна из
другой переносом и поворотом) области $A_1,\dots,A_n\subset\R^2$,
разбивающие плоскость на максимально возможное ($2^n$) число частей, которые
различаются между собой состояниями принадлежности их точек областям
$A_1,\dots,A_n$. (Заметим, что для $n\le3$ такими областями, например,
являются круги одинакового радиуса.)

\subsection*{Решения}
\begin{solution}[1 (\S\,2)]
Нумеруем конституентные области трех кругов, как на рис.~3,~\emph{б},
и записываем, какие из них образуют левую и правую части. Удобнее
обозначать область тройкой $\alpha\beta\gamma$ состояний принадлежности
к $A$, $B$, $C$ ($1$ "--- внутри, $0$ "--- вне); тогда
\[
1=111,\ 2=110,\ 3=011,\ 4=101,\ 5=010,\ 6=001,\ 7=100,\ 8=000 .
\]
\begin{center}\small
\begin{tabular}{lll}
\toprule
тождество & левая часть & правая часть\\
\midrule
$(A\cup C)\cap B\equiv(A\cap B)\cup(\ov A\cap B\cap C)$ & $\{1,2,3\}$ & $\{1,2\}\cup\{3\}$\\
$[A\cap(B\cup C)]\cup(\ov A\cap C)\equiv C\cup(A\cap B)$ & $\{1,2,4\}\cup\{3,6\}$ & $\{1,3,4,6\}\cup\{1,2\}$\\
$(A\cap C)\cup(B\cap C)\cup(C\setminus A)\equiv C$ & $\{1,4\}\cup\{1,3\}\cup\{3,6\}$ & $\{1,3,4,6\}$\\
$A\sd B\equiv(A\cup B)\setminus(A\cap B)$ & $\{4,7,3,5\}$ & $\{1,\dots,5,7\}\setminus\{1,2\}$\\
$[A\cap(B\sd C)]\cup(B\sd C)\equiv B\sd C$ & $\{2,4\}\cup\{2,4,5,6\}$ & $\{2,4,5,6\}$\\
\bottomrule
\end{tabular}
\end{center}
Во всех строках множества номеров совпадают, а так как три круга
в положении рис.~3,~\emph{б} реализуют все восемь логически возможных
состояний, совпадение областей на диаграмме и означает тождество
(см. также упр.~2 \S\,3 и решение упр.~7 \S\,3, где это рассуждение
обосновано строго). Четвертое и пятое тождества видны уже на диаграмме
двух кругов (пятое "--- это поглощение $(X\cap Y)\cup Y\equiv Y$).

\emph{Последнее тождество.} Поскольку $B\setminus\ov A\equiv B\cap A$,
напечатанное в книге равенство
$A\cup[(B\setminus\ov A)\sd C]\equiv(A\cap B)\sd C$ имеет левую часть
$A\cup\{1,2\}\sd\{1,3,4,6\}=\{1,2,4,7\}\cup\{2,3,4,6\}=\{1,2,3,4,6,7\}$,
а правую $\{2,3,4,6\}$: они различаются областями $1$ и $7$ (например, при
$A=I$, $B=C=\emptyset$). Значит, в книге опечатка. Естественный
верный вариант "--- с пересечением вместо объединения:
\[
A\cap[(B\setminus\ov A)\sd C]\equiv(A\cap B)\sd(A\cap C):
\]
слева $\{1,2,4,7\}\cap\{2,3,4,6\}=\{2,4\}$, справа
$\{1,2\}\sd\{1,4\}=\{2,4\}$. Это тождество выражает
дистрибутивность $\cap$ относительно $\sd$:
$A\cap(X\sd Y)\equiv(A\cap X)\sd(A\cap Y)$. (Верно также
$A\cap[(B\setminus\ov A)\sd C]\equiv(A\cap B)\sd(A\cap C)\equiv A\cap(B\sd C)$.)
\end{solution}

\begin{solution}[2 (\S\,2)]
\emph{Ответ: да; $n$ выпуклых областей могут дать все $2^n$
логически возможных частей.}

Для $n=4$ это показывает, например, известная диаграмма Венна из четырех
равных эллипсов (эллипс выпуклый). Общая конструкция для любого $n$:
возьмем на окружности $2^n$ различных точек $q_\alpha$, занумерованных
наборами $\alpha=(\alpha_1,\dots,\alpha_n)\in\{0,1\}^n$, и положим
\[
A_i=\operatorname{conv}\{q_\alpha:\ \alpha_i=1\}
\]
(выпуклая оболочка; чтобы получить область с внутренними точками, ее можно
слегка <<раздуть>>, взяв $\varepsilon$-окрестность при малом
$\varepsilon$ "--- выпуклость сохраняется). Точки на окружности находятся
в выпуклом положении: ни одна из них не лежит в выпуклой оболочке
остальных и даже отстоит от нее на положительное расстояние. Поэтому
$q_\alpha\in A_i\iff\alpha_i=1$, т.\,е. точка $q_\alpha$ принадлежит
конституенте $A_1^{\alpha_1}\cap\dots\cap A_n^{\alpha_n}$, и все $2^n$
конституент непусты (а вместе с малыми окрестностями точек $q_\alpha$,
лежащими вне всех $A_i$ с $\alpha_i=0$, "--- имеют внутренние точки).

Итак, максимальное число частей, различающихся состояниями
принадлежности, равно $2^n$ "--- больше быть не может, так как наборов
состояний всего $2^n$. (Связными эти части, вообще говоря, не будут; если
же считать связные компоненты, то их число для выпуклых областей не
ограничено никакой функцией от $n$: границы двух выпуклых многоугольников
могут пересекаться в сколь угодно большом числе точек.)
\end{solution}

\begin{solution}[3 (\S\,2)]
\emph{Ответ: да, для всякого $n$; можно даже обойтись одними
параллельными переносами.}

Выберем $n$ векторов $v_1,\dots,v_n$ и $2^n$ точек $x_\alpha$,
$\alpha\in\{0,1\}^n$, <<в общем положении>>, так чтобы все $n2^n$ точек
$x_\alpha-v_i$ были попарно различны (этого можно добиться сколь угодно
малым шевелением). Разобьем их на два конечных множества
\[
P_+=\{x_\alpha-v_i:\ \alpha_i=1\},\qquad P_-=\{x_\alpha-v_i:\ \alpha_i=0\}.
\]
Пусть $S$ "--- связная открытая область, содержащая все точки $P_+$ и не
содержащая ни одной точки $P_-$ (например, объединение малых кругов вокруг
точек $P_+$, соединенных тонкими <<коридорами>>, обходящими конечное
множество $P_-$, "--- плоскость без конечного множества точек связна).
Положим $A_i=S+v_i$. Тогда
\[
x_\alpha\in A_i\iff x_\alpha-v_i\in S\iff\alpha_i=1,
\]
т.\,е. $x_\alpha$ лежит в конституенте с номером $\alpha$, и все $2^n$
конституент непусты (и открыты в окрестности $x_\alpha$, если круги
выбрать малыми). Области $A_i$ конгруэнтны "--- они получаются друг из
друга переносом.

Для $n\le3$ подходят равные круги (рис.~3,~\emph{б}). Для $n\ge4$
равные (и вообще любые) \emph{круги} уже не годятся: $n$ окружностей
разбивают плоскость не более чем на $n^2-n+2$ связных частей (каждая новая,
$k$-я, окружность пересекает прежние не более чем в $2(k-1)$ точках
и добавляет не более $2(k-1)$ частей). Каждая непустая конституента
содержит хотя бы одну такую часть, поэтому непустых конституент не больше
$n^2-n+2$, а $n^2-n+2<2^n$ при $n\ge4$. Поэтому при $n\ge4$ приходится
брать области более сложной формы, как в построении выше. Для простых $n$
существуют даже симметричные диаграммы Венна из $n$ конгруэнтных кривых,
переходящих друг в друга поворотом (Григгс, Киллиан, Сэвидж, 2004).
\end{solution}

\section*{\S\,3. Принцип двойственности}
\addcontentsline{toc}{section}{\S\,3. Принцип двойственности}
\markright{\S\,3. Принцип двойственности}
\subsection*{Условия}
1. Существует ли самодвойственная двухместная операция $F(A,B)$, не равная
тождественно $A$ или $B$?

2. Описать множество всех самодвойственных формул вида $F(A_1,\dots,A_n)$
при любом $n\ge2$.

3. Описать множество всех логических формул, для которых
$F\cap F^*=\emptyset$. Будет ли при этом $F\cup F^*=I$?

4.\fixed{в оригинале номер упражнения <<4.>> пропущен.} Будем
предполагать, что $F(A_1,\dots,A_n)$ "--- \emph{строго $n$-местная}
логическая операция, если не существует таких целых чисел
$k_1,\dots,k_s\in[1,n]$ ($s<n$), что
$F(A_1,\dots,A_n)=\Phi(A_{k_1},\dots,A_{k_s})$, где $\Phi$ "--- $s$-местная
логическая формула. Для каких $n$ существуют строго $n$-местные
самодвойственные логические формулы? Простой перебор показывает, что для
$n=2$ таких формул нет. Для $n=3$ примером самодвойственной строго
трехместной формулы является (1.6).

5. Пусть $\mu(n)$ "--- количество самодвойственных $n$-местных логических
формул. Показать, что $\mu(n)$ "--- неубывающая функция.

6. Пусть $\tau(n)$ "--- количество самодвойственных строго $n$-местных
логических формул. Показать, что $\tau(2)=0$. Для каких еще $n$
$\tau(n)=0$?

7. Пусть $\lambda\{F(A_1,\dots,A_n)\}=\lambda(F)$ "--- количество
конституентных областей в диаграмме Эйлера для логической формулы
$F(A_1,\dots,A_n)$, а $\lambda(F^*)$ "--- количество таких же областей для
двойственной формулы. Описать класс функций, для которых
$\lambda(F)=\lambda(F^*)$.

\subsection*{Решения}
% Решения упражнений § 3 (с. 21--22).
% Основа: заметки claude/p21-*.md из проекта (см. каталог notes/).

\begin{remarkbox}{Общие инструменты для упражнений 1--7}
Во всех упражнениях формулы рассматриваются \emph{с точностью до
тождественного равенства}: две формулы, тождественно равные друг другу,
считаются одной и той же (они задают одну и ту же <<операцию над
множествами>>).

\textbf{Лемма (L) о двойственности.} Для любой формулы $F$
\begin{equation*}\tag{L}
F^*(A_1,\dots,A_n)\equiv\ov{F(\ov{A_1},\dots,\ov{A_n})}.
\end{equation*}

\emph{Доказательство} "--- полная индукция по числу $k(F)$ символов
операций $\cap$, $\cup$, ${}^{\overline{\phantom a}}$ в записи $F$.
Обозначим $\widehat F=\ov{F(\ov A)}$, где $\ov A=(\ov{A_1},\dots,\ov{A_n})$.

\emph{База} $k=0$: $F=A_i$ дает $F^*=A_i=\ov{\ov{A_i}}=\widehat F$;
$F=\emptyset$ дает $F^*=I=\ov\emptyset=\widehat F$; $F=I$ дает
$F^*=\emptyset=\ov I=\widehat F$.

\emph{Шаг.} Пусть (L) верно для всех формул с не более чем $k$ операциями
и $k(F)=k+1$. Внешняя операция $F$ "--- одна из трех.
\begin{itemize}[nosep]
\item $F=\ov G$, $k(G)=k$. Замена символов черту не трогает, поэтому
  $F^*=\ov{G^*}\equiv\ov{\ov{G(\ov A)}}=G(\ov A)$, и
  $\widehat F=\ov{\ov{G(\ov A)}}=G(\ov A)$.
\item $F=G\cap H$, $k(G)+k(H)=k$, так что $k(G),k(H)\le k$ (именно здесь
  нужна \emph{полная} индукция: подформулы могут быть намного короче).
  Тогда $F^*=G^*\cup H^*\equiv\ov{G(\ov A)}\cup\ov{H(\ov A)}
  \overset{10^\circ}{\equiv}\ov{G(\ov A)\cap H(\ov A)}=\widehat F$.
\item $F=G\cup H$ "--- так же, с правилом $11^\circ$. \qquad$\square$
\end{itemize}

\textbf{Коды точек и булева функция формулы.} Точке $x\in I$ сопоставим
\emph{код} $\alpha(x)=(\alpha_1,\dots,\alpha_n)\in\{0,1\}^n$, $\alpha_i=1$
при $x\in A_i$ и $\alpha_i=0$ иначе. Точки с одним кодом $\alpha$ образуют
\emph{конституенту}
\[
K_\alpha=A_1^{\alpha_1}\cap\dots\cap A_n^{\alpha_n},\qquad A^1=A,\ A^0=\ov A ;
\]
конституент $2^n$, они попарно не пересекаются и покрывают $I$.
Принадлежность точки множеству $F(A_1,\dots,A_n)$ зависит только от ее
кода (индукция по $k(F)$), поэтому формуле отвечает булева функция
$f\colon\{0,1\}^n\to\{0,1\}$ (таблица истинности; $\cap\to\wedge$,
$\cup\to\vee$, черта $\to$ отрицание, $\emptyset\to0$, $I\to1$) и
\begin{equation*}\tag{D}
x\in F\iff f(\alpha(x))=1,\qquad F=\bigcup_{\alpha:\,f(\alpha)=1}K_\alpha .
\end{equation*}

\emph{(i) Тождества = равенство функций.} $F\equiv G\iff f=g$. Если
$f(\alpha_0)\ne g(\alpha_0)$, возьмем $I=\{0,1\}^n$,
$A_i=\{\alpha:\alpha_i=1\}$; тогда $K_{\alpha_0}=\{\alpha_0\}$, и точка
$\alpha_0$ лежит ровно в одном из множеств $F$, $G$.

\emph{(ii) Функция двойственной формулы.} Замена $A_i\to\ov{A_i}$
переводит код каждой точки в $\bar\alpha=(1-\alpha_1,\dots,1-\alpha_n)$.
Поэтому по (L)
\[
f^*(\alpha)=\ov{f(\bar\alpha)} .
\]
Наборы $\{0,1\}^n$ разбиваются на $2^{n-1}$ пар противоположных
$\{\alpha,\bar\alpha\}$ (набор не совпадает со своим дополнением).
\end{remarkbox}

%----------------------------------------------------------------------
\begin{solution}[1 (\S\,3)]
\emph{Ответ: формально да "--- это $\ov A$ и $\ov B$, но обе они
по существу одноместны; самодвойственной операции, существенно
зависящей от обоих аргументов, нет.}

По (L) самодвойственность $F\equiv F^*$ равносильна
\begin{equation*}\tag{C}
F(\ov{A_1},\dots,\ov{A_n})\equiv\ov{F(A_1,\dots,A_n)},
\qquad\text{т.\,е.}\qquad f(\bar\alpha)=\ov{f(\alpha)}\ \ \forall\alpha .
\end{equation*}
При $n=2$ пары противоположных наборов: $\{00,11\}$ и $\{01,10\}$.
Свободно выбираются $f(00)$ и $f(01)$, остальное определено:
\begin{center}\small
\begin{tabular}{ccccl}
\toprule
$f(00)$ & $f(01)$ & $f(10)$ & $f(11)$ & формула\\ \midrule
0&0&1&1& $A$\\ 0&1&0&1& $B$\\ 1&0&1&0& $\ov B$\\ 1&1&0&0& $\ov A$\\ \bottomrule
\end{tabular}
\end{center}
Всего четыре самодвойственные двухместные формулы, и ни одна не зависит
от обоих аргументов. Это и есть <<простой перебор>>, упомянутый в упр.~4.
\end{solution}

%----------------------------------------------------------------------
\begin{solution}[2 (\S\,3)]
\textbf{Разминка.}
\begin{center}\small
\begin{tabular}{lll}
\toprule
$F$ & $F^*$ & самодвойственна?\\ \midrule
$A_k$, $\ov{A_k}$ & те же & да\\
$\emptyset$ & $I$ & нет\\
$A\cap B$ & $A\cup B$ & нет (при $A=\emptyset$, $B=I$)\\
$(A\cap B)\cup(\ov A\cap I)$ & $(A\cup B)\cap(\ov A\cup\emptyset)$ & нет: $\ov A\cup B$ и $\ov A\cap B$\\
$(A\cap B)\cup(A\cap C)\cup(B\cap C)$ & $(A\cup B)\cap(A\cup C)\cap(B\cup C)$ & да "--- (1.6)\\
\bottomrule
\end{tabular}
\end{center}
(Проверка (L) на четвертой строке: при подстановке $A\to\ov A$ черта над
$A$ в формуле дает $\ov{\ov A}=A$; дополняются \emph{буквы}, а не
литералы: $F(\ov A,\ov B)=(\ov A\cap\ov B)\cup A=A\cup\ov B$,
$\ov{F(\ov A,\ov B)}=\ov A\cap B=F^*$.)

\textbf{Критерий.} По (C) и (ii) $F$ самодвойственна $\iff$
$f(\bar\alpha)=\ov{f(\alpha)}$ при всех $\alpha$. Наглядно: если
одновременно поменять ролями <<внутри>> и <<снаружи>> у \emph{всех}
исходных множеств, то результат тоже меняет <<внутри>> и <<снаружи>>.
Например, для (1.6) <<точка входит хотя бы в два из трех множеств>> после
дополнения всех трех превращается в <<входит не более чем в одно>>, т.\,е.
в дополнение исходного множества "--- проверка в одну строку.

На языке конституент: дополнение всех аргументов переставляет
$K_\alpha\leftrightarrow K_{\bar\alpha}$, и самодвойственность означает,
что формула из каждой пары $\{K_\alpha,K_{\bar\alpha}\}$ берет \emph{ровно
одну}.

\begin{theorem*}[описание]
$F(A_1,\dots,A_n)$ самодвойственна тогда и только тогда, когда
\[
F\equiv\bigcup_{\alpha\in S}A_1^{\alpha_1}\cap\dots\cap A_n^{\alpha_n},
\]
где $S\subset\{0,1\}^n$ содержит ровно один элемент из каждой пары
$\{\alpha,\bar\alpha\}$. Разным $S$ отвечают разные формулы.
\end{theorem*}
\begin{proof}
<<Только тогда>>: $S=\{\alpha:f(\alpha)=1\}$ и (C). <<Тогда>>:
$f(\alpha)=1\iff\alpha\in S\iff\bar\alpha\notin S\iff f(\bar\alpha)=0$.
Различность "--- по (i).
\end{proof}
Следствия: самодвойственная формула занимает ровно половину диаграммы
Эйлера, $\lambda(F)=2^{n-1}$; их число $\mu(n)=2^{2^{n-1}}$.

\begin{theorem*}[порождение]
$F$ самодвойственна тогда и только тогда, когда
\[
\boxed{\;F\equiv\bigl(A_1\cap G(A_2,\dots,A_n)\bigr)\cup\bigl(\ov{A_1}\cap G^*(A_2,\dots,A_n)\bigr)\;}
\]
для некоторой \emph{произвольной} формулы $G$ от $n-1$ аргументов;
соответствие $G\mapsto F$ взаимно однозначно.
\end{theorem*}
\begin{proof}
Для любой формулы верно разложение Шеннона
$F\equiv(A_1\cap F_1)\cup(\ov{A_1}\cap F_0)$, $F_1=F(I,A_2,\dots)$,
$F_0=F(\emptyset,A_2,\dots)$. Подставим его в (C). Слева
$(\ov{A_1}\cap F_1(\ov{A'}))\cup(A_1\cap F_0(\ov{A'}))$, где
$A'=(A_2,\dots,A_n)$; справа
$(A_1\cap\ov{F_1})\cup(\ov{A_1}\cap\ov{F_0})$. Множество вида
$(A_1\cap X)\cup(\ov{A_1}\cap Y)$ однозначно определяет $X$ и $Y$ (при
$A_1=I$ и $A_1=\emptyset$), поэтому
$F_1(\ov{A'})=\ov{F_0(A')}$, т.\,е. $F_0=\ov{F_1(\ov{A'})}\overset{(L)}{=}F_1^*$.
Обратно, для $F$ указанного вида по (L) $G(\ov{A'})=\ov{G^*(A')}$,
$G^*(\ov{A'})=\ov{G(A')}$, и
$F(\ov A)=(\ov{A_1}\cap\ov{G^*})\cup(A_1\cap\ov G)=\ov{F(A)}$.
Наконец, $G=F(I,A_2,\dots,A_n)$ восстанавливается по $F$.
\end{proof}
Это второй вывод $\mu(n)=2^{2^{n-1}}$: самодвойственных формул от $n$
аргументов столько же, сколько вообще формул от $n-1$ аргумента.

\textbf{Проверка при $n=3$.} Перебирая 16 формул $G(B,C)$, получаем
$F=(A\cap G)\cup(\ov A\cap G^*)$:
\begin{center}\small
\begin{tabular}{lll}
\toprule
$G$ & $G^*$ & $F$\\ \midrule
$\emptyset$ & $I$ & $\ov A$\\
$I$ & $\emptyset$ & $A$\\
$B$, $\ov B$, $C$, $\ov C$ & те же & $B$, $\ov B$, $C$, $\ov C$\\
$B^b\cap C^c$ & $B^b\cup C^c$ & $\maj(\ov A,B^b,C^c)$\\
$B^b\cup C^c$ & $B^b\cap C^c$ & $\maj(A,B^b,C^c)$\\
$B\sd C$ & $\ov{B\sd C}$ & $\ov{A\sd B\sd C}$\\
$\ov{B\sd C}$ & $B\sd C$ & $A\sd B\sd C$\\ \bottomrule
\end{tabular}
\end{center}
Здесь $b,c\in\{0,1\}$, $\maj(X,Y,Z)=(X\cap Y)\cup(X\cap Z)\cup(Y\cap Z)$.
Итого 16 формул: 6 вырожденных ($A,\ov A,B,\ov B,C,\ov C$), 8
<<медиан>> $\maj(A^a,B^b,C^c)$ (среди них (1.6)) и 2 <<четностных>>
$A\sd B\sd C$, $\ov{A\sd B\sd C}$.

\textbf{Итог.} Самодвойственные формулы "--- в точности те, что
дополняют результат при дополнении всех аргументов:
$F(\ov{A_1},\dots,\ov{A_n})\equiv\ov{F(A_1,\dots,A_n)}$. Их можно выписать
(1) через конституенты, выбирая по одной из каждой пары
$\{K_\alpha,K_{\bar\alpha}\}$, или (2) как
$(A_1\cap G)\cup(\ov{A_1}\cap G^*)$ с произвольной $G$.

\textbf{Замечания.} (а) \emph{Замкнутость.} Если $F$ и $G_1,\dots,G_n$
самодвойственны, то и $F(G_1,\dots,G_n)$ самодвойственна; вместе
с проекциями $A_k$ это замкнутый класс (класс $S$ в теории Поста) "---
ср. \S\,4. (б) \emph{Связь с $R$-функциями.} При переходе к функциям
($A_i\leftrightarrow\{\omega_i\ge0\}$) дополнение множества "--- это смена
знака $\omega\mapsto-\omega$, и (C) превращается в
$f(-\omega_1,\dots,-\omega_n)=-f(\omega_1,\dots,\omega_n)$:
самодвойственным операциям отвечают \emph{нечетные} $R$-функции.
Например, $\tfrac12(x+y\mp|x-y|)$ ($\min$ и $\max$) двойственны, но не
самодвойственны, а медиана $\operatorname{med}(x,y,z)$ нечетна, как
и должно быть для (1.6).
\end{solution}

%----------------------------------------------------------------------
\begin{solution}[3 (\S\,3)]
\emph{Ответ: $F\cap F^*\equiv\emptyset$ $\iff$
$F(\ov{A_1},\dots,\ov{A_n})\equiv F(A_1,\dots,A_n)$ (<<четные>> формулы);
для них всегда $F^*\equiv\ov F$ и потому $F\cup F^*\equiv I$; более того,
условие $F\cup F^*\equiv I$ выделяет тот же класс.}

\emph{Перевод на язык функций.} По (i) $F\cap F^*\equiv\emptyset$
$\iff$ $f(\alpha)\wedge f^*(\alpha)=0$ при всех $\alpha$, а по (ii)
это $f(\alpha)\wedge\ov{f(\bar\alpha)}=0$:
\begin{equation*}\tag{1}
f(\alpha)=1\ \Rightarrow\ f(\bar\alpha)=1\qquad\forall\alpha .
\end{equation*}
Применяя (1) к $\bar\alpha$ вместо $\alpha$ (а $\bar{\bar\alpha}=\alpha$),
получаем обратную импликацию, т.\,е.
\begin{equation*}\tag{2}
f(\alpha)=f(\bar\alpha)\qquad\forall\alpha,
\qquad\text{или}\qquad F(\ov{A_1},\dots,\ov{A_n})\equiv F(A_1,\dots,A_n).
\end{equation*}
Обратно, (2) $\Rightarrow$ (1) очевидно. То же в множествах: $F\cap F^*=\emptyset$
значит $F^*\subseteq\ov F$, т.\,е. (по (L)) $\ov{F(\ov A)}\subseteq\ov{F(A)}$,
$F(A)\subseteq F(\ov A)$ при любых $A$; подставив $\ov A$ вместо $A$,
получаем обратное включение.

\begin{theorem*}
$F\cap F^*\equiv\emptyset$ тогда и только тогда, когда $F$ "---
объединение \emph{целых пар} конституент:
$F\equiv\bigcup_{\alpha\in T}(K_\alpha\cup K_{\bar\alpha})$, где $T$ "---
произвольное множество представителей пар. Таких формул $2^{2^{n-1}}$.
\end{theorem*}

\emph{Замкнутая форма.} При $n\ge2$ четные формулы "--- это в точности
\[
\boxed{\;F\equiv H\bigl(A_1\sd A_2,\ A_1\sd A_3,\ \dots,\ A_1\sd A_n\bigr)\;}
\]
с \emph{произвольной} формулой $H$ от $n-1$ аргумента. Действительно,
$\ov{A_1}\sd\ov{A_i}\equiv A_1\sd A_i$, так что такая $F$ четна. Обратно,
для четной $F$ положим $H(B_2,\dots,B_n)=F(\emptyset,B_2,\dots,B_n)$. Если
код точки $x$ есть $\alpha$ и $\alpha_1=0$, то $x\in A_1\sd A_i\iff x\in A_i$,
и правая часть содержит $x$ тогда же, когда $f(\alpha)=1$; если
$\alpha_1=1$, то $x\in A_1\sd A_i\iff x\notin A_i$, код равен
$(0,\bar\alpha_2,\dots,\bar\alpha_n)=\bar\alpha$, и по (2)
$f(\bar\alpha)=f(\alpha)$. Неформально: четная формула <<не видит>> самих
множеств "--- только то, \emph{совпадают или различаются} состояния точки
по отношению к разным множествам.

\emph{Второй вопрос.} Для четной $F$ по (ii) $f^*(\alpha)=\ov{f(\bar\alpha)}=\ov{f(\alpha)}$,
т.\,е. $F^*\equiv\ov F$; значит $F\cup F^*\equiv I$ и $F\cap F^*\equiv\emptyset$
одновременно. Обратно, $F\cup F^*\equiv I$ означает
$f(\alpha)\vee\ov{f(\bar\alpha)}=1$, т.\,е. $f(\bar\alpha)=1\Rightarrow f(\alpha)=1$,
и тем же приемом получаем (2). Итак, для любой $F$ равносильны:
\[
F\cap F^*\equiv\emptyset\iff F\cup F^*\equiv I\iff F^*\equiv\ov F\iff
F(\ov{A_1},\dots,\ov{A_n})\equiv F(A_1,\dots,A_n).
\]
Для произвольной формулы ни одно из первых двух условий не влечет
другое <<по общим соображениям>>: например, для $F=A\cap B$ имеем
$F\cap F^*=A\cap B\ne\emptyset$ и $F\cup F^*=A\cup B\ne I$.

\textbf{Примеры.} $n=1$: четные формулы $\emptyset$ и $I$ ($A$ и $\ov A$
не подходят: $A\cap A^*=A$). $n=2$: $\emptyset$, $A\sd B$,
$\ov{A\sd B}=(A\cap B)\cup(\ov A\cap\ov B)$, $I$ (т.\,е. $H=\emptyset$,
$B_2$, $\ov{B_2}$, $I$); например,
$(A\sd B)^*=(A\cup\ov B)\cap(\ov A\cup B)=\ov{A\sd B}$. $n=3$: 16 формул
$H(A\sd B,A\sd C)$; при $H=B_2\cap B_3$ получаем
$F=(A\sd B)\cap(A\sd C)=K_{100}\cup K_{011}$, и
$F^*=(A\cup\ov B\cup\ov C)\cap(\ov A\cup B\cup C)=\ov F$.

\textbf{Связь с упр.~2.} Самодвойственные формулы ($f(\bar\alpha)=\ov{f(\alpha)}$,
<<нечетные>>) и четные ($f(\bar\alpha)=f(\alpha)$) "--- два крайних
случая: у первых $F\cap F^*=F$, у вторых $F\cap F^*=\emptyset$. Классы
равномощны, и $F$ самодвойственна $\iff$ $F\sd A_1$ четна.
\end{solution}

%----------------------------------------------------------------------
\begin{solution}[4 (\S\,3)]
\emph{Ответ: строго $n$-местные самодвойственные формулы существуют при
$n=1$ и при всех $n\ge3$; единственное исключение "--- $n=2$.}

\emph{Переформулировка.} Будем говорить, что $F$ \emph{существенно
зависит} от $A_j$, если найдутся $A_1,\dots,A_n$ и $A_j'$ с
$F(\dots,A_j,\dots)\ne F(\dots,A_j',\dots)$. Строгая $n$-местность "---
это существенная зависимость от каждого аргумента: если $F$ не зависит
от $A_j$, то $F\equiv F(A_1,\dots,\emptyset,\dots,A_n)$ "--- формула от
$n-1$ аргумента; и наоборот.

\emph{$n=2$} "--- см. решение упр.~1: все четыре самодвойственные формулы
одноместны. (Самодвойственных формул от двух аргументов $2^{2^1}=4$, ровно
столько же, сколько одноместных самодвойственных $A,\ov A,B,\ov B$ "---
<<места не остается>>.)

\emph{$n=3$}: формула (1.6) $M(A,B,C)$ самодвойственна; при $B=I$,
$C=\emptyset$ она равна $A$, и по симметрии существенно зависит от
каждого аргумента.

\emph{Общий случай $n\ge3$.} По теореме о порождении (упр.~2) при
$G=A_2\cap\dots\cap A_n$, $G^*=A_2\cup\dots\cup A_n$ получаем
самодвойственную формулу
\begin{equation*}\tag{4}
F_n=(A_1\cap A_2\cap\dots\cap A_n)\cup\bigl(\ov{A_1}\cap(A_2\cup\dots\cup A_n)\bigr)
\end{equation*}
(<<точка лежит в $A_1$ и во всех остальных множествах, либо не лежит в
$A_1$, но лежит хотя бы в одном из остальных>>; при $n=3$ это
$M(\ov A,B,C)$). Прямая проверка: при $P=A_2\cap\dots\cap A_n\subseteq U=A_2\cup\dots\cup A_n$
\[
F_n^*=(A_1\cup U)\cap(\ov{A_1}\cup P)=(A_1\cap P)\cup(\ov{A_1}\cap U)\cup(U\cap P)
=(A_1\cap P)\cup(\ov{A_1}\cap U)=F_n,
\]
так как $U\cap P=P=(A_1\cap P)\cup(\ov{A_1}\cap P)$ поглощается.
Существенная зависимость: при $A_2=I$, $A_3=\dots=A_n=\emptyset$ (нужен
аргумент $A_3$, т.\,е. $n\ge3$) $F_n=\ov{A_1}$; при $A_1=I$ и всех
$A_i=I$, кроме $A_j$ ($j\ge2$), $F_n=A_j$. При $n=2$ та же формула
вырождается в $(A_1\cap A_2)\cup(\ov{A_1}\cap A_2)=A_2$.

\emph{$n=1$}: $A$ (и $\ov A$) самодвойственна и не равна тождественно ни
$\emptyset$, ни $I$.

\emph{Оценка подсчетом.} Нестрогая самодвойственная формула равна самодвойственной формуле от
$n-1$ аргумента, поэтому их не больше $n\,\mu(n-1)$, тогда как
$\mu(n)=\mu(n-1)^2$. При $n=3$: $16>3\cdot4=12$, при $n=4$: $256>64$;
при $n=2$ оценка дает $2\mu(1)=4=\mu(2)$ "--- запаса нет, и перебор
показывает, что строго двухместных действительно нет.
\end{solution}

%----------------------------------------------------------------------
\begin{solution}[5 (\S\,3)]
\emph{Что считается.} Различных \emph{записей} бесконечно много
($A_1$, $A_1\cap A_1$, \dots), но формул с точностью до тождества
$2^{2^n}$ "--- по числу таблиц $f$ (см. (i)). Самодвойственные среди них
и составляют $\mu(n)$. При этом <<$n$-местная>> формула не обязана
существенно зависеть от всех аргументов.

\emph{Монотонность без вычислений.} Всякая самодвойственная формула
$F(A_1,\dots,A_n)$ является самодвойственной и как формула от
$A_1,\dots,A_n,A_{n+1}$ (дополнение лишнего аргумента ни на что не
влияет), а различные формулы остаются различными. Значит,
$\mu(n)\le\mu(n+1)$.

\emph{Точное значение.} По критерию (C) таблица задается выбором значения
на одном наборе каждой из $2^{n-1}$ пар $\{\alpha,\bar\alpha\}$, выборы
независимы:
\[
\mu(n)=2^{2^{n-1}}:\qquad
\begin{array}{c|cccccc}
n&1&2&3&4&5&6\\\hline
\mu(n)&2&4&16&256&65\,536&4\,294\,967\,296
\end{array}
\]
Отсюда $\mu(n+1)=\mu(n)^2>\mu(n)$ "--- функция даже строго возрастает.
Для сравнения: всех $n$-местных формул $2^{2^n}=\mu(n)^2=\mu(n+1)$,
т.\,е. самодвойственных формул от $n$ аргументов ровно столько, сколько
вообще формул от $n-1$ аргумента.
\end{solution}

%----------------------------------------------------------------------
\begin{solution}[6 (\S\,3)]
\emph{Ответ: $\tau(2)=0$, и ни при каком другом $n\ge1$ $\tau(n)$ не
обращается в нуль: $\tau(1)=2$, $\tau(3)=10$, $\tau(4)=208$,
$\tau(5)=64\,386$, \dots}

$\tau(2)=0$ "--- решение упр.~1; $\tau(1)\ge1$ ($A$) и $\tau(n)\ge1$ при
$n\ge3$ (формула (4) из решения упр.~4).

\emph{Точная формула.} У каждой формулы есть однозначно определенное
множество $S\subseteq\{1,\dots,n\}$ существенных аргументов; подставив
$\emptyset$ вместо несущественных, получаем строго $|S|$-местную формулу
$\Phi$, которая самодвойственна вместе с $F$. Обратно, пара <<$S$ +
строго $|S|$-местная самодвойственная $\Phi$>> дает самодвойственную
$n$-местную формулу, и разные пары дают разные формулы. Поэтому
\[
\mu(n)=\sum_{k=0}^{n}\binom nk\tau(k),\qquad \tau(0)=0
\]
(формулы $\emptyset$, $I$ от нуля аргументов не самодвойственны). Отсюда
$\tau(n)=\mu(n)-\sum_{k<n}\binom nk\tau(k)$ (или, обращением,
$\tau(n)=\sum_k(-1)^{n-k}\binom nk\mu(k)$):
\[
\begin{aligned}
\tau(1)&=2,\qquad \tau(2)=4-2\cdot2=0,\qquad \tau(3)=16-3\cdot0-3\cdot2=10,\\
\tau(4)&=256-4\cdot10-6\cdot0-4\cdot2=208,\\
\tau(5)&=65\,536-5\cdot208-10\cdot10-10\cdot0-5\cdot2=64\,386 .
\end{aligned}
\]
(Для $n\le4$ значения проверены прямым перебором.) При $n=3$ десять
формул "--- восемь <<медиан>> $\maj(A^{\pm},B^{\pm},C^{\pm})$ и две
<<четности>>.

\emph{Почему именно $n=2$.} Вырожденных самодвойственных формул не
больше $n\,\mu(n-1)$, а $\mu(n)/(n\,\mu(n-1))=\mu(n-1)/n=2^{2^{n-2}}/n$.
При $n\ge3$ это больше единицы и быстро растет, а при $n=2$ равно
$\mu(1)/2=1$: двойка "--- единственное $n$, при котором возведение
в квадрат еще не обгоняет умножение на $n$ ($2^2=2\cdot2$). Доля строго
$n$-местных среди самодвойственных быстро стремится к единице:
$\tau(4)/\mu(4)=208/256$, $\tau(5)/\mu(5)\approx0{,}98$.
\end{solution}

%----------------------------------------------------------------------
\begin{solution}[7 (\S\,3)]
\emph{Ответ: $\lambda(F)=\lambda(F^*)$ $\iff$ $\lambda(F)=2^{n-1}$, т.\,е.
$F$ составлена ровно из половины конституентных областей (равносильно:
$\lambda(F)=\lambda(\ov F)$). Таких формул $\binom{2^n}{2^{n-1}}$.}

$\lambda(F)$ "--- число конституент, из которых состоит $F$ на полной
диаграмме Эйлера (все $2^n$ областей присутствуют), т.\,е. число единиц
в таблице $f$. Например, при $n=2$: $\lambda(A\cap B)=1$,
$\lambda(A\cup B)=3$, $\lambda(A)=2$, $\lambda(\emptyset)=0$,
$\lambda(I)=4$.

По (ii) $f^*(\alpha)=1\iff f(\bar\alpha)=0$: области $F^*$ "--- это
наборы, противоположные наборам, \emph{не} вошедшим в $F$. Отображение
$\alpha\mapsto\bar\alpha$ "--- перестановка, поэтому
\[
\boxed{\;\lambda(F)+\lambda(F^*)=2^n\;}
\]
(геометрически: $F^*$ "--- центральное отражение дополнения $\ov F$ на кубе
кодов). Значит, $\lambda(F)=\lambda(F^*)\iff\lambda(F)=2^{n-1}$; назовем
такие формулы \emph{уравновешенными}. Любое множество $S$ из $2^{n-1}$
наборов дает уравновешенную формулу, разные $S$ "--- разные формулы:
\[
\binom{2^n}{2^{n-1}}:\quad n=1:\ 2,\quad n=2:\ 6,\quad n=3:\ 70,\quad n=4:\ 12\,870,\ \dots
\]

\textbf{Примеры.} $n=2$: $A$, $\ov A$, $B$, $\ov B$, $A\sd B=K_{01}\cup K_{10}$,
$\ov{A\sd B}=K_{00}\cup K_{11}$. $n=3$: медиана (1.6) "--- 4 области из 8;
$A\sd B\sd C$ "--- 4 области; $A\cap(B\cup C)$ "--- 3 области "--- \emph{не}
уравновешена, ее двойственная $A\cup(B\cap C)$ занимает 5.

\textbf{Связь с упр.~2 и 3.} Самодвойственные формулы берут ровно одну
область из каждой пары $\{K_\alpha,K_{\bar\alpha}\}$, поэтому \emph{все}
уравновешены. Четные берут пары целиком, $m$ пар дают $\lambda=2m$
($0\le m\le2^{n-1}$; крайние случаи $\emptyset$ и $I$); уравновешены
лишь те, что составлены ровно из половины пар, $m=2^{n-2}$, их
$\binom{2^{n-1}}{2^{n-2}}$ (при $n=2$ "--- $A\sd B$ и $\ov{A\sd B}$,
но не $\emptyset$, $I$; при $n=3$ "--- 6 из 16). Уравновешенная формула не
обязана быть ни той, ни другой: при $n=3$, $S=\{000,111,001,010\}$,
\[
F=K_{000}\cup K_{111}\cup K_{001}\cup K_{010}=\ov{A\sd(B\cap C)}
\]
берет пару $\{000,111\}$ целиком (не самодвойственна), а пару
$\{001,110\}$ "--- наполовину (не четна), но $\lambda(F)=4=\lambda(F^*)$.
При $n=3$ из 70 уравновешенных формул 16 самодвойственных, 6 четных
и 48 <<прочих>>.
\end{solution}

\section*{\S\,4. Композиция и суперпозиция. Замкнутые множества функций}
\addcontentsline{toc}{section}{\S\,4. Композиция и суперпозиция. Замкнутые множества функций}
\markright{\S\,4. Композиция и суперпозиция. Замкнутые множества функций}
\subsection*{Условия}
1. Пусть $x*y\colon\R^2\to\R$ "--- некоторая ассоциативная операция,
$h(x)\colon\R\to\R$ "--- строго монотонная функция, $h^{-1}(x)$ "---
функция, обратная $h(x)$. Показать, что двухместная операция
$\varphi(x,y)\equiv h^{-1}[h(x)*h(y)]$ ассоциативна.

2. Построить ассоциативную, но не коммутативную двухместную операцию
$x*y\colon\R^2\to\R$.

3. Пусть $h_t(x)=(1-t)e^x+tx$, $h_t^{-1}(x)$ "--- обратная функция
$h_t(x)$. Показать, что функция $\varphi_t(x,y)=h_t^{-1}[h_t(x)h_t(y)]$
комму\-%
тативна и ассоциативна при $0\le t\le1$, а $\varphi_0(x,y)\equiv x+y$,
$\varphi_1(x,y)\equiv xy$.

4. Пусть $A$ "--- функциональное множество, $Q(A)$ "--- наименьшее
замкнутое множество, такое, что $A\subset Q(A)$ ($Q(A)$ "--- замкнутая
оболочка $A$). Показать, что $Q(A)=\M(H)$, где $H$ "--- полная по
отношению к $A$ система функций из $A$. Показать, что при
$H\cap\ov A\ne\emptyset$ будет $Q(A)\subset\M(H)$.\fixed{в оригинале <<при $H\cap\ov A\ne\emptyset$ $Q(A)\subset\M(H)$>> (пропущено связующее слово).} Привести примеры.

5. Пусть $A_i$ ($i=1,\dots,m$) "--- функциональные множества,
$H_i\subset A_i$ "--- соответствующие им полные системы. Показать, что
\[
\bigcup_{i=1}^{i=m}A_i\subset\bigcup_{i=1}^{i=m}\M(H_i)\subset\M\Bigl(\bigcup_{i=1}^{i=m}H_i\Bigr).
\]

6. Функция $\varphi(x_1,x_2)$ определяет ассоциативную двухместную
операцию, если $\varphi(x_1,\varphi(x_2,x_3))\equiv\varphi(\varphi(x_1,x_2),x_3)$.
Аналогично функция $\psi(x_1,x_2,x_3)$ определяет ассоциативную
трехместную операцию, если
\begin{multline*}
\psi(\psi(x_1,x_2,x_3),x_4,x_5)\equiv\psi(x_1,\psi(x_2,x_3,x_4),x_5)\equiv{}\\
\equiv\psi(x_1,x_2,\psi(x_3,x_4,x_5)).
\end{multline*}
Привести пример трехместной ассоциативной операции, не прибегая
к конструкции $\psi(x_1,x_2,x_3)\equiv\varphi(x_1,\varphi(x_2,x_3))$, где
$\varphi(x_1,x_2)$ "--- двухместная ассоциативная операция. Обобщить эти
рассуждения для $n$-местных операций ($n>3$).

7. Построить базисы для множеств целых рациональных функций с целыми,
рациональными и действительными коэффициентами.

\subsection*{Решения}
% Решения упражнений § 4 (с. 28--29).

\begin{solution}[1 (\S\,4)]
Чтобы выражение $h^{-1}[h(x)*h(y)]$ имело смысл, нужно, чтобы
$h(x)*h(y)$ лежало в области значений $h$; будем считать, что множество
$h(\R)$ замкнуто относительно $*$ (например, $h$ "--- биекция $\R\to\R$).
Строгая монотонность обеспечивает существование $h^{-1}$ на $h(\R)$
и тождество $h(h^{-1}(u))=u$. Тогда
\begin{align*}
\varphi(\varphi(x,y),z)&=h^{-1}\bigl[h\bigl(h^{-1}[h(x)*h(y)]\bigr)*h(z)\bigr]={}\\
&=h^{-1}\bigl[(h(x)*h(y))*h(z)\bigr],\\
\varphi(x,\varphi(y,z))&=h^{-1}\bigl[h(x)*h\bigl(h^{-1}[h(y)*h(z)]\bigr)\bigr]={}\\
&=h^{-1}\bigl[h(x)*(h(y)*h(z))\bigr],
\end{align*}
и правые части совпадают по ассоциативности $*$. (Это <<перенос
структуры>>: $h$ "--- изоморфизм операции $\varphi$ на операцию $*$; так же
переносится и коммутативность.)

\emph{Замечание.} Все ассоциативные операции $\varphi_0$--$\varphi_3$ из
текста параграфа имеют такой вид:
\begin{itemize}[nosep]
\item $\varphi_0$: $h(x)=\operatorname{arth}(x/c)$, $*=+$ (сложение
  <<быстрот>>; $\operatorname{th}(a+b)=\frac{\operatorname{th}a+\operatorname{th}b}{1+\operatorname{th}a\operatorname{th}b}$);
\item $\varphi_1$ ($\alpha\ne0$): $h(x)=1+\alpha x$, $*=\times$, так как
  $(1+\alpha x)(1+\alpha y)=1+\alpha(x+y+\alpha xy)$;
\item $\varphi_2$ (на $x\ge0$): $h(x)=x^2+a^2$, $*=+$;
\item $\varphi_3$: $h(x)=1/x-\alpha$, $*=+$, так как
  $1/\varphi_3=1/x_1+1/x_2-\alpha$.
\end{itemize}
\end{solution}

\begin{solution}[2 (\S\,4)]
Простейший пример "--- левая проекция $x*y=x$:
$(x*y)*z=x=x*(y*z)$, но $1*2=1\ne2=2*1$.

Менее тривиальный пример: $x*y=x\,|y|$. Ассоциативность:
$(x*y)*z=x|y|\,|z|$, $x*(y*z)=x\,\bigl|y|z|\bigr|=x|y||z|$. Некоммутативность:
$1*(-1)=1$, $(-1)*1=-1$. По упр.~1 из него получается целое семейство
ассоциативных некоммутативных операций $h^{-1}[h(x)\,|h(y)|]$.
\end{solution}

\begin{solution}[3 (\S\,4)]
При $0\le t\le1$ имеем $h_t'(x)=(1-t)e^x+t>0$, так что $h_t$ строго
возрастает. Область значений: при $t=0$ это $(0,\infty)$, при $0<t\le1$
"--- все $\R$ (при $x\to-\infty$ $h_t\sim tx\to-\infty$, при $x\to+\infty$
$h_t\to+\infty$). В обоих случаях она замкнута относительно умножения,
поэтому по упр.~1 (с $*=\times$) операция
$\varphi_t(x,y)=h_t^{-1}[h_t(x)h_t(y)]$ ассоциативна, а коммутативна,
так как коммутативно умножение.

Крайние значения: $h_0=e^x$, $h_0^{-1}=\ln$,
$\varphi_0=\ln(e^xe^y)=x+y$; $h_1(x)=x$, $\varphi_1=xy$. Таким образом,
семейство $\varphi_t$ непрерывно деформирует сложение в умножение,
оставаясь в классе коммутативных ассоциативных операций.

При $0<t<1$ обратная функция выражается через функцию Ламберта $W$
($we^w=s$): из $y=(1-t)e^x+tx$, полагая $x=y/t-w$, получаем
$we^w=\frac{1-t}{t}e^{y/t}$, т.\,е.
\[
h_t^{-1}(y)=\frac yt-W\Bigl(\frac{1-t}{t}\,e^{y/t}\Bigr).
\]
\end{solution}

\begin{solution}[4 (\S\,4)]
\emph{Существование $Q(A)$.} Множество всех правильных отображений
$\mathfrak F(\mathfrak X)$ замкнуто и содержит $A$, а пересечение любого
семейства замкнутых множеств замкнуто (доказательство теоремы~1 дословно
переносится на любое число множеств). Поэтому $Q(A)$ "--- пересечение всех
замкнутых множеств, содержащих $A$.

\emph{$\M(H)$ замкнуто.} Суперпозиции суперпозиций функций из $H$ сами
являются суперпозициями функций из $H$: $\M(\M(H))=\M(H)$ (пункты
а)--г) определения применяются к уже построенным суперпозициям).

\emph{Равенство $Q(A)=\M(H)$ при $H\subset A$, $A\subset\M(H)$.}
(1) $\M(H)$ замкнуто и содержит $A$, значит, $Q(A)\subset\M(H)$.
(2) $H\subset A\subset Q(A)$ и $Q(A)$ замкнуто, значит,
$\M(H)\subset Q(A)$.

\emph{Случай $H\cap\ov A\ne\emptyset$} (система полна по отношению к $A$,
но содержит функции не из $A$). Рассуждение (1) остается в силе:
$Q(A)\subset\M(H)$; рассуждение (2) "--- нет, и включение может быть
строгим.

\emph{Примеры.} $A=C^R$ (полиномы), $H=H_0=\{x_1+x_2,\ x_1x_2,\ a\}$:
$H\subset A$, $Q(C^R)=\M(H_0)=C^R$ (множество $C^R$ само замкнуто).
$A=\{x_1+x_2,\ x_1x_2\}\cup\R$ (константы): $Q(A)=C^R$ "--- замкнутая
оболочка конечного набора операций может быть гораздо шире самого набора.
$A=C^R$, $H=H_0\cup\{\sin x\}$: $H\cap\ov A=\{\sin x\}$, и
$Q(A)=C^R\subsetneq\M(H)\ni\sin x$.
\end{solution}

\begin{solution}[5 (\S\,4)]
Первое включение "--- определение полноты: $A_i\subset\M(H_i)$ для
каждого $i$. Второе следует из монотонности операции $\M$: если
$H'\subset H''$, то всякая суперпозиция функций из $H'$ есть суперпозиция
функций из $H''$, т.\,е. $\M(H')\subset\M(H'')$; применяем это
к $H_i\subset\bigcup_jH_j$.

Включения, вообще, строгие: при $H_1=\{x_1+x_2\}$, $H_2=\{x_1x_2\}$
функция $x_1x_2+x_3$ принадлежит $\M(H_1\cup H_2)$, но не принадлежит
$\M(H_1)\cup\M(H_2)$.
\end{solution}

\begin{solution}[6 (\S\,4)]
\emph{Пример.} $\psi(x_1,x_2,x_3)=x_1-x_2+x_3$ на $\R$. Все три способа
расстановки дают одно и то же:
\begin{multline*}
\psi(\psi(x_1,x_2,x_3),x_4,x_5)=\psi(x_1,\psi(x_2,x_3,x_4),x_5)={}\\
=\psi(x_1,x_2,\psi(x_3,x_4,x_5))=x_1-x_2+x_3-x_4+x_5 .
\end{multline*}
(В любой группе аналогично $\psi=x_1x_2^{-1}x_3$; это <<груда>> Прюфера
и Бэра.)

\emph{Эта операция не имеет вида $\varphi(x_1,\varphi(x_2,x_3))$
с ассоциативной $\varphi$.} Допустим противное:
$\varphi(x,\varphi(y,z))=\varphi(\varphi(x,y),z)=x-y+z$. Пусть $R$ "---
множество значений $\varphi$. Если $u=\varphi(y,z)\in R$, то
$\varphi(x,u)=x+(z-y)$ при всех $x$; значит, $\varphi(x,u)=x+k(u)$ для
$u\in R$ (величина $k(u)=z-y$ определена однозначно). Аналогично
$\varphi(v,z)=z+m(v)$ для $v\in R$. При $u,v\in R$ отсюда
$v+k(u)=\varphi(v,u)=u+m(v)$, т.\,е. $u-k(u)=v-m(v)=c$ "--- постоянная, и
$\varphi(x,u)=x+u-c$ при $u\in R$. Тогда
$x-y+z=\varphi(x,\varphi(y,z))=x+\varphi(y,z)-c$, т.\,е.
$\varphi(y,z)=z-y+c$ при всех $y,z$; следовательно, $R=\R$ и одновременно
$\varphi(x,u)=x+u-c$ и $\varphi(x,u)=u-x+c$ при всех $x,u$, что
невозможно. Противоречие.

\emph{Обобщение на $n$-местные операции.} Операция $\psi$ от $n$
аргументов \emph{ассоциативна}, если при подстановке $\psi$ в любое из $n$
мест получается одна и та же функция $2n-1$ аргументов:
\begin{multline*}
\psi(\psi(x_1,\dots,x_n),x_{n+1},\dots,x_{2n-1})\equiv{}\\
\equiv\psi(x_1,\psi(x_2,\dots,x_{n+1}),\dots,x_{2n-1})\equiv\dots\\
\dots\equiv\psi(x_1,\dots,x_{n-1},\psi(x_n,\dots,x_{2n-1})).
\end{multline*}
Тривиальные примеры дают <<итерации>> двухместной ассоциативной операции,
$\psi=x_1*x_2*\dots*x_n$. Нетривиальные: пусть $\zeta$ "--- корень
степени $n-1$ из единицы ($\zeta^{n-1}=1$). Тогда
\[
\psi(x_1,\dots,x_n)=x_1+\zeta x_2+\zeta^2x_3+\dots+\zeta^{n-1}x_n
\]
ассоциативна: при подстановке внутреннего $\psi$ на $k$-е место его
$j$-й аргумент получает множитель $\zeta^{k-1}\zeta^{j-1}=\zeta^{(k+j-1)-1}$
"--- ровно тот, что положен месту $k+j-1$, а аргументы, сдвинувшиеся на
$n-1$ мест вправо, сохраняют множитель, так как $\zeta^{n-1}=1$.
В действительных числах можно взять $\zeta=-1$ при нечетном $n$:
$\psi=x_1-x_2+x_3-\dots+x_n$ (при $n=3$ "--- пример выше). Вообще, если
$h$ "--- биекция, то $h^{-1}[\psi(h(x_1),\dots,h(x_n))]$ снова
ассоциативна (как в упр.~1). Систематически $n$-местные ассоциативные
операции изучаются в теории $n$-арных групп (Дёрнте, 1928; Пост, 1940).
\end{solution}

\begin{solution}[7 (\S\,4)]
Функции из базиса должны принадлежать самому множеству $\M_0$, а для
проверки того, что система $H$ "--- базис, достаточно убедиться, что она
полна и что после удаления \emph{любого одного} элемента полнота
теряется. Заметим главное: при суперпозиции полиномов коэффициенты
результата лежат в кольце, порожденном (сложением и умножением)
коэффициентами исходных полиномов.

\textbf{Целые коэффициенты} ($\Z[x_1,x_2,\dots]$):
\[
H=\{x_1+x_2,\ x_1x_2,\ -1\}.
\]
Полнота: $1=(-1)(-1)$, $0=1+(-1)$, $n=1+\dots+1$, $-n=(-1)\,n$; далее
одночлены "--- произведениями, многочлены "--- суммами. Минимальность:
без $-1$ нет ни одной константы (из $x_1+x_2$ и $x_1x_2$ без констант
получаются только многочлены с неотрицательными коэффициентами без
свободного члена); без $x_1x_2$ получаются лишь функции вида
$c+\sum n_ix_i$ степени $\le1$; без $x_1+x_2$ "--- лишь одночлены
$\pm x_1^{k_1}\cdots x_s^{k_s}$ и константы $\pm1$.

\textbf{Рациональные коэффициенты} ($\Q[x_1,x_2,\dots]$):
\[
H=\{x_1+x_2,\ x_1x_2,\ -1\}\cup\{1/p:\ p\ \text{простое}\}.
\]
Полнота: кольцо, порожденное $-1$ и всеми $1/p$, есть $\Q$
($1/q$ для составного $q$ "--- произведение $1/p$). Минимальность: без
$1/p$ все коэффициенты лежат в кольце $\Z[1/q:\ q\ne p]$, где у
знаменателей нет множителя $p$, так что $1/p$ не получить; без $-1$ все
коэффициенты неотрицательны; без $x_1x_2$ или $x_1+x_2$ "--- как выше.
\emph{Конечного} базиса здесь нет: у конечной системы конечное число
коэффициентов, их знаменатели делят некоторое $N$, и все суперпозиции
имеют коэффициенты из $\Z[1/N]\ne\Q$.

\textbf{Действительные коэффициенты} ($C^R$). Система
$H_0=\{x_1+x_2,\ x_1x_2,\ a,\ \forall a\in\R\}$ полна (пример~1), но
базисом в смысле определения не является: отдельные константы избыточны
(например, $2=1+1$). Если, как в обозначении $H_0$ книги, считать семейство
всех констант одним элементом системы, то $H_0$ "--- базис: без сложения,
без умножения или без констант полноты нет (рассуждения те же, что выше).
Если же требовать неизбыточности поэлементно, то по замечанию о
коэффициентах базис вида $\{x_1+x_2,\ x_1x_2\}\cup C$ существует тогда
и только тогда, когда у кольца $\R$ есть неизбыточная система
образующих $C$ (такая, что ни одно $c\in C$ не выражается через остальные
с помощью сложения и умножения). Для $\Z$ и $\Q$ такие системы указаны
выше ($\{-1\}$ и $\{-1\}\cup\{1/p\}$); для $\R$ этот теоретико-кольцевой
вопрос выходит за рамки книги, и мы оставляем его открытым (в предисловии
автор предупреждает, что часть задач носит проблемный характер).
\end{solution}


\appendix
\chapter{Замечания}
\section*{Замечания к исходным рабочим заметкам}
Решения упражнений \S\,3 основаны на заметках из каталога
\texttt{notes/claude/} (файлы \texttt{p21-*.md}). При переносе исправлено:
\begin{itemize}
\item В разминочной таблице заметки к упр.~2 формула
  $(A\cap B)\cup(\ov A\cap I)$ отмечена как самодвойственная (<<обе
  $\equiv\ov A\cup B$>>). Это неверно: $F\equiv\ov A\cup B$, а
  $F^*=(A\cup B)\cap\ov A\equiv\ov A\cap B$ (что, кстати, правильно
  вычислено ниже в той же заметке при проверке леммы (L)). В книге эта пара
  приведена лишь как пример \emph{двойственных} формул.
\item В решении упр.~4 добавлена прямая проверка самодвойственности
  формулы $F_n$ и замечание о вырождении при $n=2$; в упр.~6 "--- явная
  ссылка на конструкцию $F_n$; в остальном решения следуют заметкам.
\end{itemize}


\end{document}
